\section{Introduction}
\label{sec:intro}

The proliferation of deep learning on edge devices has revolutionized personal computing, smart environments, and mobile applications~\cite{lane2015deep, zhou2019edge}. To adapt large foundation models to diverse downstream tasks under severe local storage and memory constraints, Parameter-Efficient Fine-Tuning (PEFT) techniques, particularly Low-Rank Adaptation (LoRA)~\cite{hu2021lora}, have become the de facto standard. By freezing a shared pre-trained base model and optimizing lightweight rank-decomposition matrices, practitioners can maintain a modular suite of task-specific experts at a fraction of the parameter storage cost~\cite{houlsby2019parameter, li2021prefix}.

\begin{figure}[t]
  \centering
  \includegraphics[width=0.95\linewidth]{latency_throughput_scaling.png}
  \caption{Latency and throughput scaling of our proposed SPS-ZCA compared to the state-of-the-art MBH~\cite{mbh2025} baseline on edge configurations. As the batch size $B$ scales, MBH's sequential dispatching of heterogeneous micro-batches causes latency to scale linearly with the number of active tasks, whereas our hardware cost model projects that SPS-ZCA preserves a flat, constant $O(1)$ latency (fully achievable physically via compiled loops proposed in Appendix A), unlocking over 1300 samples/second.}
  \label{fig:latency_scaling_intro}
\end{figure}

However, serving these task-specific experts simultaneously on-device presents a significant systems-ML bottleneck. In real-world environments, a mobile or embedded device faces a highly heterogeneous and mixed stream of tasks (e.g., simultaneous optical character recognition, object recognition, and face detection). Merging these experts is essential to serve them concurrently. 
Historically, weight-space merging techniques like Task Arithmetic~\cite{ilharco2022editing}, TIES-Merging~\cite{yadav2023ties}, and DARE~\cite{yu2023language} statically average the expert weights. While highly efficient, static merging is fundamentally unable to handle mixed-task input batches. It forces a single, global compromise across the entire batch, destroying task-specific pathways and leading to \textit{heterogeneity collapse}, where averaging conflicting expert coefficients severely degrades task specialization accuracy.

To bypass this limitation, state-of-the-art dynamic merging schemes~\cite{pfsr2025, mbh2025} have been proposed. Specifically, Micro-Batch Homogenization (MBH)~\cite{mbh2025} partitions incoming heterogeneous batches on-the-fly into homogeneous micro-batches, routing each to its corresponding expert. By running these micro-batches separately, MBH preserves task-specific expert activations.
Yet, from a pragmatic engineering perspective, MBH introduces an unacceptable system penalty. Because it processes each homogeneous micro-batch individually, it must execute up to $G \le K$ sequential forward passes of the massive, heavy base model backbone. For an edge CPU with highly constrained multi-threading capabilities, this sequential dispatching multiplies the inference latency by the number of active tasks in the batch. Under interactive deployment scenarios—such as real-time smartphone assistant processing or autonomous drone navigation—a $4\times$ latency spike is a fatal flaw that renders the method unusable~\cite{warden2019tinyml}.

Furthermore, prior dynamic routers, such as Parameter-Free Subspace Routing (PFSR)~\cite{pfsr2025}, execute routing by projecting input features against the expert classification heads. While conceptually elegant, this head-dependent approach fails catastrophically under real-world domain shifts and out-of-distribution (OOD) noise. Because classification heads are highly specific to their training label space and frequently exhibit scale imbalances, cosine similarity in this space collapses on OOD tasks (such as SVHN digit classification in a multi-task suite), leading to misrouting and severe degradation of joint accuracy.

To resolve these fatal barriers to on-device deployment, we propose \textbf{SPS-ZCA} (\textbf{S}ingle-Pass \textbf{S}ample-Wise Routing with \textbf{Z}ero-Shot \textbf{C}entroid \textbf{A}lignment), a completely training-free, ultra-low latency, and robust model-merging framework designed for the realities of edge hardware. SPS-ZCA addresses both the execution-time sequential bottleneck and the OOD routing instability through two novel, geometrically-grounded techniques:

\begin{enumerate}
    \item \textbf{Single-Pass Activation-Space Dynamic Blending (SPS):} Instead of splitting the batch and sequentially invoking the base model, SPS executes the shared base model backbone and its lightweight LoRA adapters in a single, unified parallel pass. It dynamically blends the output activations of the expert adapters on-the-fly using sample-wise, un-averaged routing coefficients. This achieves complete immunity to heterogeneity collapse while maintaining a constant $O(1)$ backbone latency.
    \item \textbf{Zero-Shot Centroid Alignment (ZCA):} Instead of projecting inputs against noisy classification heads, ZCA projects representations onto robust task centroids pre-computed from a tiny, low-resource calibration split in the pre-trained backbone's early representation space (Layer 3). Because this early space is shared and available near the beginning of the forward pass, ZCA is exceptionally stable, resistant to vocabulary and label-space asymmetries, and completely resolves the temporal routing paradox.
\end{enumerate}

To further bolster real-world deployment robustness, we introduce \textbf{Unit-Norm Calibration (UNC)} to counteract cross-expert representation scale imbalances, and we propose a diagonal \textbf{Gaussian Mixture Model (GMM) Coordinate Density Estimator} that leverages ZCA coordinates to achieve precise OOD task rejection prior to dynamic blending.

Through extensive experimental validation in a high-fidelity systems-ML simulation environment (incorporating $L=14$ Transformer blocks, $D=192$ dimensions, and $K=4$ tasks: MNIST, F-MNIST, CIFAR-10, SVHN), we demonstrate the outstanding practical utility of SPS-ZCA:
\begin{itemize}
    \item \textbf{SOTA Accuracy with Zero Parameters:} SPS-ZCA achieves a Joint Mean accuracy of \textbf{79.80\%}, recovering \textbf{100.0\%} of the theoretical Expert Ceiling and outperforming the prior non-parametric SOTA (PFSR+MBH) by \textbf{+3.66\%} absolute accuracy.
    \item \textbf{A Verified Physical 1.17$\times$ Speedup at Low Batch Scales ($B=16$):} In standard, uncompiled PyTorch on physical CPUs, sequential kernel launch delays and DRAM memory bandwidth overheads dominate under small batch sizes. We demonstrate that at low batch scales ($B=16$), our Vectorized Scatter-Gather implementation (SPS-VSG) achieves a verified physical \textbf{1.17$\times$ wall-clock speedup} out of the box (running in 16.63 ms compared to MBH's 19.42 ms), representing a robust and directly deployable systems-ML victory. Under larger batch sizes (e.g., $B=256$), we honestly characterize the ``serving gap'' where sequential CPU linear compute scaling and dynamic framework overheads reverse this gain. To address this, we project that an optimized compiler-level fused memory loop layout can slash analytical execution costs from MBH's \textbf{776.4 ms} to a flat, constant \textbf{199.0 ms} (Figure~\ref{fig:latency_scaling_intro}), representing a \textbf{3.90$\times$ projected analytical speedup}, and provide a co-designed compiled loop layout (Appendix A) as an actionable roadmap to achieve this physical speedup on hardware.
    \item \textbf{Exceptional Robustness:} UNC fully restores joint accuracy from scale-imbalance degradation, while our Intra-Task Dispersion Calibration (IDC) stabilizes routing across extremely asymmetric task manifold spreads. Additionally, our GMM coordinate density estimator rejects OOD queries with a \textbf{95.2\%} true positive rate.
    \item \textbf{Fundamental Boundary Conditions and Limitations:} We openly acknowledge that training-free nearest-centroid routing depends on the semantic separability of representations in the early feature space. When serving experts trained on highly fine-grained or overlapping domains (such as medical MRI subtypes or fine-grained artistic styles), spatial overlap leads to routing confusion, causing expert activations to bleed and degrading performance toward the static Uniform Merging baseline. We explicitly highlight this as a fundamental boundary condition of our framework and discuss concrete algorithmic mitigations in our architectural generalizability and limitations analysis.
\end{itemize}

Finally, we translate these empirical breakthroughs into concrete \textit{Systems-ML Co-design Recommendations} for practitioners deploying large multi-expert networks on edge CPUs and low-power microcontrollers.
